% Copyright (c) 2024 by Daniel R. Driver.
% !TEX encoding = UTF-8 Unicode
% !TEX TS-program = XeLaTeX

\documentclass[titlepage]{article}

% This document presumes a file structure and set of inputs that are
% available at: git@github.com:danieldriver/syllabi.git

\newcommand\policy{../policy}
\newcommand\incl{../includes}
\ProvidesFile{variables.tex}[2022/09/11 v2.3 -- Syllabus variables]

\usepackage{xspace} % make manual spaces (like \mycmd\ ) unnecessary
\usepackage{xifthen} % provides \isempty test

% variables for internal use
\newcommand\prof{}
\newcommand\pdegree{}
\newcommand\pphone{}
\newcommand\pemail{}
\newcommand\poffice{}
\newcommand\phours{}
%
\newcommand\ccode{}
\newcommand\ctitle{}
\newcommand\cseries{}
\newcommand\cversion{}
\newcommand\csemester{}
\newcommand\cmeetson{}
\newcommand\cmeetsat{}
\newcommand\cmeetsin{}
\newcommand\cwebsite{}
\newcommand\cdescrip{}
\newcommand\cprereqs{}
\newcommand\edobject{}

% in case of fully online, BTh, or undergraduate courses - https://tex.stackexchange.com/a/5896
\newif\ifonline
\newcommand\Int[2]{\ifonline#1\else#2\fi}
\newif\ifbth
\newcommand\bth[2]{\ifbth#1\else#2\fi}
\newif\ifUG
\newcommand\UG[2]{\ifUG#1\else#2\fi}

% commands for setting variables in the preamble
\newcommand\professor[2][PhD]{
  \renewcommand\pdegree{#1\xspace}
  \renewcommand\prof{#2\xspace}}
\newcommand\phone[1]{
  \renewcommand\pphone{\addfontfeatures{Numbers=Monospaced}#1\xspace}}
\newcommand\email[1]{
  \renewcommand\pemail{\href{mailto:#1}{#1}\xspace}}
\newcommand\officehours[2][AST Library/Teams]{
  \renewcommand\poffice{#1\xspace}
  \renewcommand\phours{\href{https://outlook.office.com/bookwithme/user/b97b147a2e5c43b4a0e200c2cb630b77@astheology.ns.ca?anonymous}{#2}\xspace}%
  }
%
\newcommand\coursecode[2][1.0]{
  \renewcommand\cversion{#1\xspace}
  \renewcommand\ccode{#2\xspace}}
\newcommand\coursetitle[2][]{
  \ifthenelse{\isempty{#1}}%
    {}% do nothing if #1 is empty, else:
    {\renewcommand\cseries{#1\\[1ex]}}
  \renewcommand\ctitle{#2\xspace}}
\newcommand\semester[1]{
  \renewcommand\csemester{#1\xspace}}
\newcommand\meets[3]{
  \newcommand\AM{\textsc{am}}
  \newcommand\PM{\textsc{pm}}
  \renewcommand\cmeetson{#1\xspace}
  \renewcommand\cmeetsat{\Int{From 9:00 \AM}{#2}\xspace}
  \renewcommand\cmeetsin{\Int{\href{https://smu.brightspace.com/d2l/login}{Brightspace}}{#3}\xspace}}
\newcommand\website[2][]{
  \renewcommand\cwebsite{
  \ifthenelse{\isempty{#1}}%
    {\href{https://#2}}% copy url if #1 is empty, else
    {\href{#1}}% use specified url and print #2
  {#2}\xspace}}
\newcommand\cdescription[2][BF 1001; RM 1000 or GTRS 6000]{
  \renewcommand\cprereqs{#1}
  \renewcommand\cdescrip{#2\par}}
\newcommand\objectives[1]{
  \renewcommand\edobject{#1\par}}


%\onlinetrue % \Int{true}{false}
\coursecode[2.1.1]{HB 3114}
\coursetitle{The Twelve Prophets}

% Taught as RLGS 3143 in:
%   - Winter 2010
%   - Winter 2013
% Taught as HB 3114 in:
%   - Winter 2018
%   - Winter 2024

\professor{Daniel R. Driver}
\phone{902-425-7051}
\email{daniel.driver@astheology.ns.ca}
\officehours{Thursdays, 1:00--3:00 \PM}

\semester{Winter Term 2024}
\meets{Thursdays}% \meets{on}{at}{in}
      {4:00--6:30 \PM}
      {Classroom 3/Teams}
\website[https://danieldriver.com/courses/hb-3114/]{danieldriver.com}
\cdescription{% copy from the current Academic Calendar
	To better understand the origins and impact of the book of the
	Twelve Prophets, this course will look at its redactional history
	and editorial shaping, its final canonical forms, and its effects as
	biblical prophecy. It will introduce current issues in scholarly
	debate about the Twelve while also exploring ways that they meet the
	ancient expectation that their “bones” can “send forth new life”
	amidst the people of God (Sirach 49:10).
}% end of course description
%\objectives{% recall Bloom's taxonomy: http://www.celt.iastate.edu/teaching/RevisedBlooms1.html
%By the end of the course students should be able to:
%}% end of learning objectives

\ProvidesFile{preamble.tex}[2014/09/02 v1.1 -- Syllabus preamble]

% basic typography
\usepackage{fontspec,realscripts}
\defaultfontfeatures{Ligatures=TeX}
\setmainfont[%
  Numbers={Proportional,OldStyle},
  SmallCapsFeatures={RawFeature=+c2sc;+smcp,LetterSpace=1}]{Scala OT}
\newfontfamily\lining[Numbers={Lining}]{Scala OT}
\setsansfont[%
  SmallCapsFeatures={RawFeature=+c2sc;+smcp,LetterSpace=1}]{Scala Sans Pro}
\setmonofont[Scale=MatchLowercase]{Menlo}
\usepackage{sectsty}
\allsectionsfont{\sffamily\scshape}
\frenchspacing
\setlength{\emergencystretch}{3em} % prevent overfull lines

% custom font size and leading
\renewcommand\tiny{\fontsize{6}{9}\selectfont}
\renewcommand\scriptsize{\fontsize{7}{10}\selectfont}
\renewcommand\footnotesize{\fontsize{8}{10}\selectfont}
\renewcommand\small{\fontsize{9}{11}\selectfont}
\renewcommand\normalsize{\fontsize{10}{12}\selectfont}% base size
\renewcommand\large{\fontsize{11}{13}\selectfont}
\renewcommand\Large{\fontsize{12}{14}\selectfont}
\renewcommand\LARGE{\fontsize{16}{18}\selectfont}% "course syllabus \\ semester" benefits from more lead
\renewcommand\huge{\fontsize{20}{22}\selectfont}
\renewcommand\Huge{\fontsize{24}{26}\selectfont}

% layout packages: page, logo, tables
\usepackage[scale={0.6,0.8},
            xetex]{geometry}
\usepackage{graphicx}
\usepackage{array}     % allow insertions of column styling with >{}
\usepackage{booktabs}  % elegant horizontal rules in tables

% custom macros for a session count in the schedule of readings
\newcounter{session}
\newcounter{columns}
\newcounter{spancolumns}
\newcounter{midcolumns}
\newcounter{courseunit}
\newcommand\setcolumncount[2][0]{ % optionally set count to other than 0,
  \setcounter{session}{#1}        % e.g. to -1, or to a standing count
  \setcounter{columns}{#2}
  \setcounter{spancolumns}{#2-1}
  \setcounter{midcolumns}{#2-2}}
\newcommand\sessioncount{\stepcounter{session}\arabic{session}}
\newcommand\sessionskip[1]{\multicolumn{1}{@{}r@{ }}{#1}}
\newcommand\unit[1]{\multicolumn{\thecolumns}{c}{%
  \scshape\stepcounter{courseunit}\roman{courseunit}. #1}}
\newcommand\noclass[1]{\sessionskip{---.}&\multicolumn{\themidcolumns}{@{}l}{\itshape No Class: #1}}
\newcommand\rarr{\Large\char"EA31\hspace*{0.5pt}} % pointer finger w/ reminder tie in Scala OT
\newcommand\reminder[2]{\sessionskip{\rarr}&\multicolumn{\thespancolumns}{@{}l}{#1 \hfill #2}}

% color to match Tyndale's branding
\usepackage[usenames]{xcolor}
% predefined: black, white, red, green, blue, cyan, magenta, yellow
\definecolor{TyndaleURLs}{HTML}{0062A0} % links on tyndale.ca
\definecolor{TyndaleBlue}{cmyk}{1,1,0,.32}
\definecolor{TyndaleGold}{cmyk}{0,.27,1,0}
\definecolor{TyndaleRed}{cmyk}{0,1,.99,.04}
\definecolor{TyndaleBlack}{cmyk}{0,0,0,1}
\definecolor{TyndaleGreen}{cmyk}{.45,0,1,.24}
\definecolor{TyndaleOrange}{cmyk}{0,.79,1,0}
\definecolor{TyndaleAqua}{cmyk}{.47,0,.24,0}
\definecolor{TyndaleYellow}{cmyk}{.03,.03,.35,0}

% metadata (assumes a host of definitions are made in the main file)
\usepackage[setpagesize=false,     % leave this to geometry
            hyperfootnotes=false,  % fragile and distracting
            xetex]{hyperref}
\hypersetup{breaklinks=true,       % allow link text to break across lines
            colorlinks=true,       % colorlinks resets pdfborder to 0 0 0
            urlcolor=TyndaleURLs,  % for external links
            linkcolor=TyndaleRed,  % for normal internal links
            citecolor=TyndaleGold, % for bibliographical citations in text
            pdfauthor={\prof},
            pdftitle={\ccode: \ctitle},
            pdfsubject={Tyndale UC, \csemester},
            pdfcreator={github.com/danieldriver/syllabus}}
\urlstyle{same}                    % don't use monospace font for urls

% custom footlines
\usepackage{fancyhdr}
\pagestyle{fancy} % turn it on
\fancyhf{}        % reset everything
\renewcommand{\headrulewidth}{0pt} % remove header line as well
\lfoot{\sffamily\scshape\footnotesize\ctitle, \csemester}
\rfoot{\sffamily\scshape\footnotesize\prof\quad\thepage}

% language support: polyglossia needs to go last
\usepackage{polyglossia}
\setdefaultlanguage{english}
\setotherlanguage{hebrew}
\newfontfamily\hebrewfont[Script=Hebrew]{SBL BibLit}

% protect positioning of margin table in policy/grades -- needs to come after polyglossia
\usepackage{marginfix}

% gratuitous with custom title page, but useful as a fallback
\title{\ccode: \ctitle}
\author{\professor}
\date{\semester}


\begin{document}
\ProvidesFile{title.tex}[2013/09/06 v1.0 -- Syllabus title page]

\begin{titlepage}
  \begin{center}

    \begin{minipage}{\textwidth}
      \parbox[t]{0.5\textwidth}{
        \mbox{}\\[-\baselineskip] % dummy line to align parboxes
        \includegraphics[width=0.5\textwidth]{.syllabus/includes/TyndaleUC}}
      \hfill
      \parbox[t]{0.4\textwidth}{
        \raggedleft\LARGE\sffamily Course Syllabus\\
        \semester}
    \end{minipage}

    \vfill

    {\sffamily\textsc{\LARGE\MakeLowercase\coursecode}\\[2ex]
      \bfseries\itshape\Huge\coursetitle}

    \vfill

    \begin{tabular}{>{\bfseries}rl>{\bfseries}rl}
      \toprule
      Instructor & \professor                   & Course  & \liningnums{\coursecode}~\texttt{(v\version)} \\
      \midrule
      Phone      & \texttt{\phone}              & Meets   & \meeton                                       \\
      Email      & \href{mailto:\email}{\email} & Time    & \meetat                                       \\
      Office     & \officenumber                & Room    & \meetin                                       \\
      Hours      & \officehours                 & Website & \href{http://\website}{\website}              \\
      \bottomrule
    \end{tabular}

    \vfill

    \begin{description}
      \item[Commuter Hotline]
        Class cancellations due to inclement weather or illness will
        be announced on the commuter hotline at \texttt{416.226.6620
        x2187}. Alternately, weather cancellation information is posted
        at \href{http://tyndale.ca/weather}{tyndale.ca/weather}.
      \item[MyTyndale.ca]
        This course may have materials stored on its website, such as
        handouts or readings that may be needed in order to complete
        assignments. Students are responsible for checking these course
        pages on a regular basis. Here, too, students are able to view
        their grades throughout the semester. For more information see
        Section~\ref{mytyndale}, below.
      \item[Mail]
        Students are responsible for information communicated through
        their campus mailboxes and student e-mail accounts. A mailbox
        directory hangs beside the mailboxes. For more information
        contact the Registrar's office.
    \end{description}

  \end{center}

  \section{Course Description}
  \label{description}

  \emph{From the Academic Calendar:} \cdescription

\end{titlepage}
\setcounter{page}{2} % count the title page as page 1


%  \section{Learning Objectives}
%  \label{objectives}
%
%  \edobject

\section{Required Texts \& Materials}
\label{texts}

In addition to \textbf{the Bible in a modern translation} such as the
NJPS or NRSV, the following texts are required. Students are strongly
encouraged to purchase their own copies. AST Library copies will be
placed on a two-hour reserve.

\begingroup
\renewcommand{\section}[2]{}% temporarily remove the section heading
\begin{thebibliography}{Theodoret}% use the longest item in the bibliography

	\bibitem[Heschel]{Heschel} Heschel, Abraham J.
	\emph{The Prophets}.
	New York: Harper \& Row, 1962. Repr., New York: Perennial Classics, 2001.
	\textsc{isbn} 978-0060936990. \href{https://amzn.to/46qhye3}{Order online.}

	\bibitem[Theodoret]{Theodoret} Theodoret of Cyrus.
	\emph{Commentary on the Twelve Prophets}. Vol. 3 of \emph{Commentaries on the Prophets}.
	Translated by Robert C. Hill.
	Brookline, MA: Holy Cross Orthodox Press, 2006.
	\textsc{isbn} 978-1885652768. \href{https://amzn.to/47kYajz}{Order online.}

%	\bibitem[Two Sides]{Two} Ehud Ben Zvi and James D. Nogalski.
%	\emph{Two Sides of a Coin: Juxtaposing Views on Interpreting the Book of the Twelve / the Twelve Prophetic Books}.
%	Piscataway, NJ: Gorgias Press, 2009.
%	\textsc{isbn} 978-1607243038.

%	\bibitem[Seitz]{Seitz} Christopher R. Seitz.
%	\emph{The Goodly Fellowship of the Prophets: The Achievement of Association in Canon Formation}.
%	Grand Rapids: Baker Academic, 2009.
%	\textsc{isbn} 978-0801038839.

%	\bibitem[Jerome]{Jerome1} Jerome; Thomas P. Scheck, trans.
%	\emph{Commentaries on the Twelve Prophets, Vol. 1 (Ancient Christian Texts)}.
%	Downers Grove: IVP Academic, 2016.
%	\textsc{isbn} 978-0830829163.

\end{thebibliography}
\endgroup

\section{Supplementary Texts}
\label{supplementary}

Supplementary readings will be recommended throughout the semester.
Excerpts from this literature may be placed on reserve or made available
for download through the course website. Outstanding students will also
pursue leads on their own, using library databases and pursuing notes
and bibliographies in the required reading.

The following titles are notable in relation to the main texts. They are
not required.

\begingroup
\renewcommand{\section}[2]{}% temporarily remove the section heading
\begin{thebibliography}{3}% use the longest item in the bibliography

	\bibitem{Bickerman} Bickerman, Elias.
	\emph{Four Strange Books of the Bible: Jonah, Daniel, Koheleth, Esther}.
	New York: Schocken, 1967.

	\bibitem{Sabbath} Heschel, Abraham Joshua.
	\emph{The Sabbath: Its Meaning for Modern Man}.
	New York: Farrar, Straus and Giroux, 1951.

	\bibitem{Hill} Hill, Robert C.
	\emph{Reading the Old Testament in Antioch}. Bible in Ancient Christianity 5.
	Leiden: Brill, 2005. Link: \href{https://ast.primo.exlibrisgroup.com/view/action/uresolver.do?operation=resolveService&package_service_id=785748290007188&institutionId=7188&customerId=7185&VE=true}{AST eBook}.

	\bibitem{Nogalski-Sweeney} Nogalski, James D., and Marvin A. Sweeney, eds.
	\emph{Reading and Hearing the Book of the Twelve}. SBL Symposium Series.
	Atlanta: SBL, 2000.

	\bibitem{Redditt-Schart} Redditt, Paul L., and Aaron Schart, eds.
	\emph{Thematic Threads in the Book of the Twelve}.
	Berlin: Walter de Gruyter, 2003.

	\bibitem{Watts-House} Watts, James W., and Paul R. House, eds.
	\emph{Forming Prophetic Literature: Essays on Isaiah and the Twelve in Honor of John D.W. Watts}.
	Sheffield: Sheffield Academic, 1996.

\end{thebibliography}
\endgroup

Also, the following basic works are also worth consulting and even owning.
\cite{rlgs} in particular contains sound advice on core skills like
reading religious texts, writing essays and reviews, revising essays,
making oral presentations, and learning languages.

\begingroup
\renewcommand{\section}[2]{}% temporarily remove the section heading
\begin{thebibliography}{Making}% use the longest item in the bibliography

	\bibitem[Making]{rlgs} Northey, Margot, Bradford A. Anderson, and Joel N. Lohr.
	\emph{Making Sense in Religious Studies: A Student's Guide to Research and Writing}.
	3rd ed. Toronto: Oxford University Press, 2019. AST Library: Reference BL 41 N67 2019.

	\bibitem[SBL2]{sbl2} Collins, Billie Jean, et al.
	\emph{The SBL Handbook of Style}.
	2nd ed. Atlanta: SBL Press, 2014.
	Designed to augment \href{https://proxy.openathens.net/login?qurl=https%3A%2F%2Fwww.chicagomanualofstyle.org%2Fbook%2Fed17%2Ffrontmatter%2Ftoc.html}{\emph{Chicago Style}}
	(the standard at AST), there is also a free
	\href{https://www.sbl-site.org/assets/pdfs/pubs/SBLHSsupp2015-02.pdf}{Student Supplement for SBL2}. AST Library: Reference PN 147 S26 2014.

\end{thebibliography}
\endgroup

\section{Course Outline}
\label{outline}

We will adhere to the schedule in \autoref{schedule} as closely as
possible, though the professor reserves the right to adjust it to suit
the needs of the class.

\setcolumncount{4}% set up \sessioncount, \unit{}, \noclass{}, and \reminder{memo}{date} macros
\begin{table}[htbp]% set to `p' to put the schedule on its own page
  \centering
  \begin{tabular}{>{\sessioncount.}r@{ }llr}% make sure the column config agrees with \setcolumncount
	\toprule
	\sessionskip{\textbf{\S}.}&\textbf{Primary}&\textbf{Secondary}&\textbf{Date}\\
	\midrule
		& None: Introductions & Syllabus & 11 Jan. \\
		& Hosea 1             & Theodoret 1--42, Heschel 47--61   & 18 Jan. \\
		& Hosea 2--3          & Theodoret 42--48, Heschel 61--75  & 25 Jan. \\
		& Hosea 4--6          & Theodoret 48--57, Heschel 3--19   &  1 Feb. \\
		& Hosea 7--9          & Theodoret 57--67, Heschel 19--31  &  8 Feb. \\
		& Hosea 10--14        & Theodoret 67--83, Heschel         & 15 Feb. \\
	\reminder{Essay 1: Due by the end of Week Six}{18 Feb.} \\
	\noclass{Reading Week (Tuesday to Friday)}                    & 22 Feb. \\
		& Joel 1              & Theodoret 85--90, Heschel         & 29 Feb. \\
		& Joel 2--end         & Theodoret 90--102, Heschel        &  7 Mar. \\
	\noclass{Attend Grad Projects in lieu of regular class}       & 14 Mar. \\
		& Amos 1--2           & Theodoret 103--10, Heschel 32--39 & 21 Mar. \\
		& Amos 3--4           & Theodoret 110--16, Heschel 39--46 & 28 Mar. \\
		& Amos 5--6           & Theodoret 116--22, Heschel        &  4 Apr. \\
	\reminder{Essay 2: Due by the end of Week Eleven}{7 Apr.} \\
		& Amos 7--9           & Theodoret 122--28, Heschel        & 11 Apr. \\
	\reminder{End of Term: Final marks are due for all courses}{17 Apr.} \\
	\bottomrule
  \end{tabular}
  \caption{Schedule of Readings}
  \label{schedule}
\end{table}

See the AST website for a list of other \href{http://www.astheology.ns.ca/students/academic-dates.html}{important dates}.

\section{Evaluation}
\label{evaluation}

\subsection{Grade Structure for \ccode}
\label{structure}

\begin{enumerate}

	\item Weekly \textbf{seminar discussions} will facilitate close work
	with text and interpretation in select books from the Twelve.
	Students should come to class fully prepared to discuss the assigned
	material, which will ordinarily include one to four chapters of
	biblical material per week, plus some commentary on those chapters.
	Commentary available for study and preparation is by no means
	limited to the assigned works of \cite{Theodoret} and
	\cite{Heschel}. Also, Heschel does not follow the text sequentially;
	his book should be read beyond the assigned pages in a self-directed
	manner.

	\item At some point in the semester, each student should read and
	report back to the class about some recent scholarship on the book
	of the Twelve. This \textbf{reading report} should be scheduled
	between weeks two and twelve with only one student per week.
	Presentations will be allotted 20 minutes. The report itself should
	last approximately 15 minutes, allowing 5 minutes for questions
	after. The reading should be drawn from the literature in
	\autoref{supplementary} or \autoref{bib}, selected in consultation
	with the professor. 

	\item A \textbf{first essay} will be due at the end of the sixth
	week of class (Sunday, 18 February 2024). It should be 4,000 words
	long, plus or minus 250 words. It should answer the following
	question: \emph{How does the book of Hosea speak to the goodness of
	God and the sovereignty of God?} The paper must advance a thesis.

	\item A \textbf{second essay} will be due at the end of the eleventh
	week of class (Sunday, 7 April 2024). It should be 4,000 words long,
	plus or minus 250 words. It should answer the following question:
	\emph{How does the presentation of God develop in the first part of
	the book of the Twelve?} Note that this question is not a repeat of
	the first one. Rather, it builds on the first paper by examining the
	character and presentation of God in the literary context that
	follows Hosea. If this question does not appeal, students are
	invited to propose an alternate question to the professor. The
	professor reserves the right to turn down the proposal or to make a
	counterproposal. The paper must advance a thesis.
	
\end{enumerate}

The breakdown for the semester's total work is shown in
\autoref{grade-dist}.

\begin{table}[htbp]
  \centering
  {\lining
  \begin{tabular}{lr}
    \toprule
    Weekly Seminars & 10\% \\
    Reading Report  & 10\% \\
    Essay 1         & 40\% \\
    Essay 2         & 40\% \\
    \bottomrule
  \end{tabular}}
  \caption{Distribution of Grades}
  \label{grade-dist}
\end{table}

\ProvidesFile{grades.tex}[2023/09/06 v2.5 -- Course policy]

AST's \href{https://www.astheology.ns.ca/students/resources.html}{Academic
Calendar} provides guidelines and criteria for academic assessment. Marks are
assigned by letter grade using these benchmarks. Note that graduate and
undergraduate programs are currently assessed on a different scales at AST (see
\autoref{grade-syst}). For graduate students, a grade of C is deemed
``Unsatisfactory,'' whereas for undergraduate (BTh) students, a C is
``Satisfactory'' and the grade of D is possible for ``Marginal'' performance.
Other grades have different definitions, too.

\begin{table}[htbp]
 \centering
 {\lining\footnotesize
 \begin{tabular}{lll|ll}
   \toprule
               & \multicolumn{2}{c|}{Graduate Courses}   & \multicolumn{2}{c}{Undergraduate Courses} \\
               \cline{2-5}
   A+          & 94--100    & Exceptional               & 90--100    & Excellent                 \\
   A           & 87--93     & Outstanding               & 85--89     & Excellent                 \\
   A\char"2212 & 80--86     & Excellent                 & 80--84     & Excellent                 \\ [1ex]
   B+          & 77--79     & Good                      & 75--79     & Good                      \\
   B           & 73--76     & Acceptable                & 70--74     & Good                      \\
   B\char"2212 & 70--72     & Marginally Acceptable     &            &                           \\ [1ex]
   C           & 60--69     & Unsatisfactory            & 60--69     & Satisfactory              \\
   D           &            &                           & 50--59     & Marginal                  \\
   F           & 0--59      & Failure                   & 0--49      & Failure                   \\
   FP          & 0          & Failure due to Plagiarism & 0          & Failure due to Plagiarism \\
   \bottomrule
 \end{tabular}}
 \caption{Outline of Grades at AST}
 \label{grade-syst}
\end{table}

\ProvidesFile{other.tex}[2018/0/03 v2.4 -- Course policy]

\section{Other Course Policy}
\label{policy}

Late work will not be accepted, except in genuinely extenuating
circumstances. Students must submit something before the deadline if
they wish to receive credit. Unless I state otherwise, written
assignments are to be uploaded by 11:59 \PM\ on the date indicated.

Essay submissions must be typewritten, double-spaced, and formatted as
PDFs. They should be free from error. In this course they should follow
either \emph{The SBL Handbook of Style} (refer to the online
\href{https://www.sbl-site.org/assets/pdfs/pubs/SBLHSsupp2015-02.pdf}{
Student Supplement}), or \emph{The Chicago Manual of Style}, per AST's
policy for major papers. As a reminder, AST also upholds an Inclusive
Language Policy.

\href{http://www.eerdmans.com/Pages/Item/59043/Commentary-Statement.aspx}{Plagiarism},
if \href{https://www.theguardian.com/world/2013/feb/09/german-education-minister-quits-phd-plagiarism}{detected},
will result in failure of the course.

Students should request permission to record a class or lecture. If
permission is granted, or if recordings are provided (as in the case of
an online course), I stipulate that all recordings be for personal use
only. They may not be shared or distributed.

If you have abilities or disabilities that require modifications to the
assessment process or other aspects of this course, please advise the
course instructor as soon as possible.

Finally, I encourage the conscientious use of laptops, tablets, and
other technology in my classes. In classroom settings, realize that, as
\href{http://dx.doi.org/10.1016/j.compedu.2012.10.003}{cognitive
psychologists have demonstrated}, ``laptop multitasking hinders
classroom learning for both users and nearby peers.'' Do your part to
foster an environment of open dialogue by honouring the presence of your
classmates. In online settings, consider both the physical environment
in which you choose to work and the virtual environment that you help
create through your participation in various forums. Let your engagement
in this course be marked by rigour and charity alike.

\section{Further Bibliography}
\label{bib}

Access and study the following resources:

Schart, Aaron. “Bibliography on the Book of the Twelve Prophets.” This database with nearly all literature on the book of the Twelve since 2010 is maintained and shared freely by the author at \url{https://www.zotero.org/groups/twelveprophets/items}.

Sweeney, Marvin A. “Book of the Twelve Prophets.” In \emph{Oxford Bibliographies: Biblical Studies}, \href{https://go.openathens.net/redirector/astheology.ns.ca?url=https://www.oxfordbibliographies.com/view/document/obo-9780195393361/obo-9780195393361-0016.xml}{https://www.oxfordbibliographies.com/view/document/obo-9780195393361/obo-9780195393361-0016.xml} (last modified 30 March 2015). AST students can sign in with their OpenAthens credentials.

\end{document}
